\documentclass[11pt,a4paper]{article}
\usepackage[utf8]{inputenc}
\usepackage[vietnamese]{babel}
\usepackage{amsmath,amssymb}
\usepackage{geometry}
\geometry{a4paper, total={180mm,260mm}, left=15mm, right=15mm, top=18mm, bottom=18mm}

\usepackage{tikz}
\usetikzlibrary{shapes.geometric, arrows.meta, positioning, calc}
\tikzset{
  arr/.style={-{Stealth[scale=0.8]}},
  every node/.append style={align=center}
}
\usepackage{tcolorbox}
\tcbuselibrary{skins,breakable}
\usepackage{fontawesome}
\usepackage{fancyhdr}
\usepackage{enumitem}

\pagestyle{fancy}
\fancyhf{}
\renewcommand{\headrulewidth}{0.5pt}
\renewcommand{\footrulewidth}{0.5pt}
\fancyhead[L]{\small\bfseries\color{violet!85!black}\faBook\ SỔ TAY KIẾN THỨC TOÁN 2}
\fancyhead[R]{\small\bfseries\color{violet!85!black}Lớp Toán Thầy Lĩnh -- loptoanthaylinh.vn}
\fancyfoot[C]{\small\color{gray}-- Sổ tay kiến thức: Bài 24. Luyện tập chung Chủ đề 4 --}

\setlength{\parindent}{0pt}
\setlength{\parskip}{0pt}

\begin{document}

% --- BANNER BÀI HỌC ---
\begin{tcolorbox}[
  enhanced,
  colback=violet!8!white,
  colframe=violet!70!black,
  arc=3mm, boxrule=1.2pt,
  left=3mm, right=3mm, top=2mm, bottom=2mm,
  halign=center
]
  {\LARGE\bfseries\color{violet!90!black}\faBookmark\ BÀI 24: LUYỆN TẬP CHUNG CHỦ ĐỀ 4}\\[1mm]
  {\itshape\footnotesize\color{black!75} Đại tổng kết kĩ năng cộng, trừ có nhớ trong phạm vi 100 và giải toán lời văn trên 1 trang A4}
\end{tcolorbox}

\vspace{-2mm}

\noindent
\begin{minipage}[t]{0.49\textwidth}\vspace{0pt}
\begin{tcolorbox}[
  enhanced,
  colback=violet!4!white,
  colframe=violet!75!black,
  arc=2mm, boxrule=0.8pt,
  left=2mm, right=2mm, top=1.8mm, bottom=1.8mm,
  title={\footnotesize\bfseries\color{white}\faKey\ 1. Hai kĩ năng trụ cột}
]
{\footnotesize
\begin{itemize}\setlength{\itemsep}{1pt}\setlength{\leftskip}{-2.5mm}
  \item \textbf{Phép cộng có nhớ:}\\
  - Hàng đơn vị cộng lại $\ge 10 \implies$ nhớ $1$ chục.\\
  - Hàng chục: cộng thêm $1$ nhớ vào kết quả.
  \item \textbf{Phép trừ có nhớ:}\\
  - Hàng đơn vị không trừ được $\implies$ mượn $1$ chục.\\
  - Hàng chục: thêm $1$ vào số trừ rồi mới trừ đi (hoặc bớt 1 ở số bị trừ).
  \item Luôn tính theo thứ tự từ \textbf{phải sang trái}.
\end{itemize}
}
\end{tcolorbox}
\end{minipage}
\hfill
\begin{minipage}[t]{0.49\textwidth}\vspace{0pt}
\begin{tcolorbox}[
  enhanced,
  colback=teal!4!white,
  colframe=teal!75!black,
  arc=2mm, boxrule=0.8pt,
  left=2mm, right=2mm, top=1.8mm, bottom=1.8mm,
  title={\footnotesize\bfseries\color{white}\faSliders\ 2. Kĩ năng giải toán tổng hợp}
]
{\footnotesize
\begin{itemize}\setlength{\itemsep}{1pt}\setlength{\leftskip}{-2.5mm}
  \item \textbf{Tính giá trị dãy tính 3 số:}\\
  Thực hiện lần lượt từ trái qua phải:\\
  $54 + 28 - 36 = 82 - 36 = 46$.
  \item \textbf{Tìm thành phần chưa biết:}\\
  $\square - 25 = 48 \implies \square = 48 + 25 = 73$.\\
  $64 - \square = 27 \implies \square = 64 - 27 = 37$.
  \item \textbf{Toán lời văn nhiều hơn, ít hơn.}
\end{itemize}
}
\end{tcolorbox}
\end{minipage}

\vspace{-1.5mm}
\begin{tcolorbox}[
  enhanced,
  colback=orange!4!white,
  colframe=orange!80!black,
  arc=2mm, boxrule=0.8pt,
  left=2mm, right=2mm, top=1.5mm, bottom=1.5mm,
  title={\footnotesize\bfseries\color{white}\faSitemap\ 3. Sơ đồ tư duy đối sánh: Cơ chế Nhớ trong Phép cộng và Phép trừ}
]
\centering
\begin{tikzpicture}[x=1.5cm, y=0.8cm]
  \node[rectangle, rounded corners=2pt, draw=violet!80!black, fill=violet!15, text width=4.2cm, align=center, font=\scriptsize\bfseries] (cong) at (0,0) {
    \textbf{PHÉP CỘNG CÓ NHỚ}\\Hàng đơn vị đầy chục ($\ge 10$)\\$\implies$ \textbf{Thêm 1} vào hàng chục
  };
  \node[rectangle, rounded corners=2pt, draw=teal!80!black, fill=teal!15, text width=4.2cm, align=center, font=\scriptsize\bfseries] (tru) at (4.5,0) {
    \textbf{PHÉP TRỪ CÓ NHỚ}\\Hàng đơn vị thiếu (không trừ được)\\$\implies$ \textbf{Mượn 1} chục để trừ
  };
  \draw[<->, very thick, orange!80!black] (cong) -- node[above, font=\tiny\bfseries] {Kiểm tra chéo} node[below, font=\tiny\bfseries] {Thử lại kết quả} (tru);
\end{tikzpicture}
\end{tcolorbox}

\vspace{-1.5mm}
\noindent
\begin{minipage}[t]{0.49\textwidth}\vspace{0pt}
\begin{tcolorbox}[
  enhanced,
  colback=green!4!white,
  colframe=green!60!black,
  arc=2mm, boxrule=0.8pt,
  left=2mm, right=2mm, top=1.8mm, bottom=1.8mm,
  title={\footnotesize\bfseries\color{white}\faCheckCircle\ 4. Ví dụ thực hành chuẩn mực}
]
{\footnotesize
\textbf{Bài toán thực tế:}\\
Đàn gà nhà Mai có $54$ con, trong đó có $28$ con gà mái, còn lại là gà trống. Hỏi nhà Mai có bao nhiêu con gà trống?\\
\textit{Bài giải:}\\
Nhà Mai có số con gà trống là:
$$54 - 28 = 26\text{ (con)}$$
Đáp số: $26$ con gà trống.
}
\end{tcolorbox}
\end{minipage}
\hfill
\begin{minipage}[t]{0.49\textwidth}\vspace{0pt}
\begin{tcolorbox}[
  enhanced,
  colback=red!4!white,
  colframe=red!75!black,
  arc=2mm, boxrule=0.8pt,
  left=2mm, right=2mm, top=1.8mm, bottom=1.8mm,
  title={\footnotesize\bfseries\color{white}\faExclamationTriangle\ 5. Mẹo tránh bẫy trắc nghiệm}
]
{\footnotesize
\begin{itemize}\setlength{\itemsep}{1pt}\setlength{\leftskip}{-2.5mm}
  \item \textbf{Bẫy tính nhầm khi có số 0 ở số bị trừ:}\\
  $60 - 24 = 36$ ($0$ không trừ được $4$, lấy $10 - 4 = 6$, nhớ $1$; $6 - 3 = 3$). Tránh lấy $4 - 0 = 4$!
  \item Luôn đặt tính ra giấy nháp cẩn thận từng hàng.
\end{itemize}
}
\end{tcolorbox}
\end{minipage}


\end{document}
